\section{拟映射越墙（Yang Zhou）}%
\label{sec:Quasimap wall-crossing (Yang Zhou)}

本讲中的主要越墙（wall-crossing）定理与主空间构造来自~\cite{Zhou2022}。更早的亏格零与高亏格结果见~\cite{CFK2014,CFK2020}，稳定拟映射（quasimap）的基础则见~\cite{CFKM2014}。

设 \(X\) 为五次三维簇。我们的目标是计数 \(X\) 上的亏格零曲线，这一问题已由 Givental 与 Lian-Liu-Yau 解决。通过闭浸入 \(X \subset \P^4\)，我们可以写出
\[ \bar{M}_0(X, d) \subset \bar{M}_0(\P^4, d), \]
它是光滑的，并带有 \((\C^{\times})^5\) 的环面作用。若一个映射在该作用下不动，则其定义域必然有一些分支映到不动点，另一些分支则构成 \(1\) 维轨道的 \(n :1\) 覆盖。特别地，不动点轨迹由一组图来标号。

在这一情形下，我们有量子 Lefschetz 原理，它断言
\[ [\bar{M}_0(X, d)]^{\on{vir}} = e(E) \cap [\bar{M}_0(\P^4, d)]. \]
这里向量丛 \(E\) 由
\[ E = \pi_* f^* \mc{O}(5), \]
定义，其中 \(\pi\) 是泛曲线，\(f\) 是泛稳定映射。我们将对下述结果给出另一种视角：若
\begin{align*}
    I(q, z) &= 1 + \sum_{d > 0} \frac{\prod_{k=1}^{5d} (5H+kz)}{\prod_{k=1}^d (H+kz)^5} q^d \in H^*(X)[z^{-1}] \ps{q} \\
    &= I_0(q) \cdot 1 + I_1(q) \frac{H}{z} + I_2(q) \frac{H^2}{z^2} + I_3(q) \frac{H^3}{z^3},
\end{align*}
则
\[ J (Q, z) = 1+ \sum_{d>0} \phi^{\alpha} \ab<\frac{\phi_{\alpha}}{z(z-\psi)}>_{0,1}^X Q^d = \frac{I}{I_0} \]
其中镜像变换为 \(Q = q \exp\ab(\frac{I_1}{I_0})\)。

\subsection{朴素计数}%
\label{sub:Naive counting}

我们将考虑从固定 \(\P^1\) 出发的“映射”\(\P^1 \to X \subset \P^4\)。为简单起见，设 \(X = \ab(\sum x_i^5 = 0)\) 为 Fermat 五次三维簇。它由一族 \([f_0, \ldots, f_4]\) 给出，其中 \(f_i \in \C[u,v]\) 满足 \(\deg f_i = d\) 且 \(\sum f_i^5 = 0\)。为了得到真正的态射，我们需要 \(f_i\) 没有公共零点，但这样得到的模空间不是真的（proper），因此我们放宽这一条件。于是我们考虑模空间
\[ QG(X, d) =  \ab\{ [f_0, \ldots, f_4] \middle\vert \Vectorstack{{ f_i \in \C[u,v] } { \deg f_i = d }}, \sum f_i^5 = 0, (f_0, \ldots, f_4) \neq (0,\ldots, 0) \} \bigg/ \C^{\times}. \]
它含于相应的模空间
\begin{align*}
    QG(\P^4, d) &\coloneqq  \ab\{ [f_0, \ldots, f_4] \mid f_i \in \C[u,v], \deg f_i = d, (f_0, \ldots, f_4) \neq (0, \ldots, 0) \} \\
    &\cong \P^{5d+5-1} 
\end{align*}
之中，后者参数化到射影空间的有理映射。这里有两个问题。其一，定义域也被参数化了，因此我们需要考虑 \(\P^1\) 的自同构。其二，为了得到稳定映射，我们需要对定义域作爆破。

首先，令 \(\C^{\times}\) 通过缩放 \(u\) 作用于 \(\P^1\)。考虑由下式定义的特殊不动点轨迹 \(F_{\star} \subset QG(X, d)\)：
\[ F_{\star,d} = \ab\{ [f_0, \ldots, f_4] \text{ fixed and } (v=0) \text{ is not a common zero} \}. \]
于是，\(F_{\star,d} = \ab\{ [a_0 u^d, \ldots, a_4 u^d] \mid [a_0, \ldots, a_4] \in X \} \cong X\)。由此，我们可以利用形变理论中的标准论证把 I-函数改写为
\begin{align*}
    I(q,z) &= \sum_{d > 0} \ab(\frac{[F_{\star, d}]^{\vir}}{e_{\C^{\times}}(N^{\vir})})^{\mr{PD}} q^d \\
    &= 1 + \sum_{d > 0} \ab( \frac{[X]}{e_{\C^{\times}}(N^{\vir})})^{\mr{PD}} q^d
\end{align*}
它位于 \(H^*_{\C^{\times}}(X) \ps{q}[z^{-1}]\) 中。

\subsection{不严格的越墙}%
\label{sub:Wall-crossing without rigor}


其次，考虑一条带标记点的结点曲线 \(C\)，外加一条带有 \(\C^{\times}\) 作用的尾，并设我们有映射 \(f \colon C \to X\)。若取极限
\[ \lim_{t \to 0} [f_0(t^{-1}u, v), \ldots, f_4(t^{-1}u, v)] = [a_0 u^{d_0}, \ldots, a_4 u^{d_0}], \]
则可见气泡的极限落在 \(F_{\star, d_0}\) 中，其中公共零点远离结点。于是我们可以用一个新的标记点替换这条尾，该标记点携带 I-函数的某些项。对五次三维簇而言，这恰好是 \(I_1 H - I_0 \psi = [zI-z]_{+, z = -\psi}\)。这一过程结束时，我们剩下 \(C = \P^1\) 以及 I-函数的各种贡献。

注意，若考虑 \(\C^{\times}\) 作用于 \(\P^4\)，则 \(f \colon C \to \P^4\) 是不动的，只要对所有 \(t \in \C^{\times}\) 存在 \(\varphi_t\) 使得下图交换：
\begin{equation*}
\begin{tikzcd}
    C \ar{r}{f} \ar{d}[swap]{\varphi_t} & \P^4 \\
    C \ar{ur}[swap]{t \cdot f}
\end{tikzcd}
\end{equation*}
于是我们在定义域上得到 \(\C^{\times}\) 的作用。

我们将定义到 \(X\) 的 \(\ep\)-稳定拟映射的模空间 \(Q_{g,n}^{\ep}(X, d)\)，它类似于稳定映射模空间。我们希望把这些模空间看作彼此虚双有理的，其中每个都是带完美障碍理论的真 DM叠，因而我们可以定义拟映射不变量
\[ \ab< \gamma_1 \psi_1^{a_1} \cdots \gamma_n \psi_n^{a_n}>_{g,n}^{\ep} \coloneqq \int_{[Q_{g,n}^{\ep}(X,d)]^{\vir}} \prod_{i=1}^n \ev_i^*(\gamma_i) \psi_i^{a_i}. \]
这里 \(\ep \in \Q_{>0}\) 是正有理数。可能的 \(\ep\) 所成的空间具有墙与室的结构，墙出现在 \(\frac{1}{n}\)（\(n\) 为正整数）处。当 \(\ep > 1\) 时我们得到格罗莫夫–威滕（GW）理论（即稳定映射），而随着向左移动，我们允许的有理尾变少，但基点变多。最终在 \(\ep \to 0^+\) 的极限下，不允许任何有理尾。最后，我们在 \(\frac{1}{d_0}\) 处的墙附近得到如下越墙公式：
\[ \ab<->_{g,n,d}^{\ep_+} - \ab<->_{g,n,d}^{\ep_-} = \sum_{k} \frac{1}{k!} \ab<-, \mu, \ldots, \mu>_{g,n+k, d-kd_0}^{\ep_+}. \]
其中 \(\mu = [zI-z]_{+, q^{d_0}, z=-\psi}\)。在 \(g = 0\)、\(n = 1\)、\(\ep = 0^+\) 的特殊情形下，我们得到
\[ \sum \phi^{\alpha} \ab<\frac{\phi_{\alpha}}{z(z-\psi)}>^{\ep}_{0,1} = I_{z^{<0}}. \]
从 GW 室出发越过所有墙，我们得到
\[ I = 1 + \sum_{\substack{d \geq 1 \\ k \geq 0}} \frac{1}{k!} \sum_{\alpha} \phi^{\alpha} \ab< \frac{\phi_{\alpha}}{z(z-\psi)}, I_1 H - I_0 \psi, \ldots, I_1 H - I_0 \psi>_{0,1+k,d}^{\infty} q^d, \]
它通过膨胀子方程与除子方程给出镜像变换。

\subsection{拟映射}%
\label{sub:Quasimaps}

\begin{defn}
    到五次三维簇 \(X\) 的一个\textit{预稳定拟映射}由
    \[(C, x_1, \ldots, x_n, \mc{L}, \varphi_0, \ldots, \varphi_4),\] 
    给出，其中 \((C, x_1, \ldots, x_n)\) 是预稳定曲线，\(\mc{L} \in \Pic C\)，\(\varphi_0, \ldots, \varphi_4 \in H^0(C, \mc{L})\)，且满足
    \[ \varphi^{-1}(0) \coloneqq \ab\{x \in C \mid \varphi_0(x) = \cdots = \varphi_4(x) = 0 \} \]
    有限并远离特殊点。此外，我们要求 \(\sum \varphi_i^5 = 0\)。
\end{defn}

重要的是，我们把 \(\varphi^{-1}(0)\) 称为\textit{基点}的集合。在 \(\varphi^{-1}(0)\) 之外，我们有真正的态射 \(C \setminus \varphi^{-1}(0) \to X\)。从几何上看，假设我们有一个定义在 \(\A^1 \setminus 0\) 上的族，当 \(t \to 0\) 取极限时，零轨迹 \(\varphi_i^{-1}(0)\) 最终会碰撞在一起。我们既可以在中心纤维中保留这个拟映射，也可以对中心纤维作若干次爆破以得到真正的态射，代价是冒出一条有理尾。因此，任何预稳定拟映射的模空间都不是分离的，所以我们必须考虑某种稳定性条件。

\begin{defn}
    一个预稳定拟映射 \((C, \mathbf{x}, \mc{L}, \varphi)\) 称为\textit{\(\ep\)-稳定的}，如果
    \begin{enumerate}
        \item 对所有 \(x \in \varphi^{-1}(0)\)，有 \(\on{len}(x) \leq \frac{1}{\ep}\)，其中 \(\on{len}(x) = \min \ab\{ \on{ord}_x (\varphi_i)\}\)；
        \item \(\Q\)-线丛 \(\mc{L}^{\ep} \otimes \omega_C^{\log}\) 在 \(C\) 的每个不可约分支上都有正次数。
    \end{enumerate}
\end{defn}

注意，对每个不可约分支 \(C' \subseteq C\)，有 \(\deg \mc{L}|_{\mc{C}'} \geq 0\)。因此发生变化的是带 \(\leq 2\) 个特殊点的亏格 \(0\) 曲线，或不带标记点的亏格 \(1\) 曲线。注意在 \((0,0)\) 与 \((1,0)\) 情形下有 \(C = C'\)，而在 \((0,2)\) 情形下稳定性等价于 \(\deg \mc{L} > 0\)，故在 \((0,1)\) 情形下第二个条件等价于
\[ \deg \mc{L}|_{C'} > \frac{1}{\ep}. \]

\begin{exm}
    在 \(\ep > 1\) 的室中，注意第一个条件蕴含没有基点，第二个条件蕴含所有不稳定分支都有正次数，因此我们就回到了稳定映射。

    若移动到 \(\frac{1}{2} < \ep < 1\) 的室，则必须收缩次数 \(1\) 的有理尾。例如，若有一条由两条次数 \(1\) 的有理曲线组成、以一个尾结束的链，我们就收缩最外侧的尾，剩下一条次数 \(2\) 的有理尾，并带有一个长度 \(1\) 的基点。

    若移动到 \(\frac{1}{3} < \ep < \frac{1}{2}\) 的室，则必须收缩所有次数 \(2\) 的有理尾，并用长度 \(2\) 的基点替换它们。
\end{exm}

\begin{rmk}
    在射影空间的情形，存在真正的态射
    \[ Q_g^{\ep_+}(\P^N, d) \to Q_g^{\ep_-}(\P^N, d) \qquad (C_g \cup C_0, \mc{L}, \varphi) \mapsto (C_g, \mc{L}|_{C_g}(p), \varphi(p)). \]
    对一般靶空间这是不可能的，但这一想法为拟映射越墙现象提供了直觉。
\end{rmk}

现在我们讨论如何把它推广到更一般的靶空间。设
\[ X = V \sslash_{\theta} G, \]
其中 \(V\) 仿射，\(G\) 约化，\(\theta\) 是 \(G\) 的特征标。我们需要要求
\begin{enumerate}
    \item \(V^{\mr{ss}}(\theta) = V^s(\theta) \neq \emptyset\)；
    \item \(V\) 至多具有 lci 奇点，且 \(V^s(\theta)\) 光滑。
\end{enumerate}
注意我们要求稳定轨迹具有有限稳定化子。为方便起见，我们还要求
\begin{enumerate}
    \item \(V \sslash G\) 是射影的；
    \item \(G\) 在稳定轨迹上自由作用。
\end{enumerate}

\begin{exm}
    在五次的情形，\(V\) 是五次三维簇的仿射锥。不稳定轨迹由原点组成，它是唯一的奇点。
\end{exm}

\begin{exm}
    另一个例子是考虑带一般稳定性条件的 \(\C^N \sslash (\C^{\times})^r\)，得到一个环面簇。
\end{exm}

\begin{exm}
    我们可以考虑格拉斯曼簇
    \[ G(2,4) = M_{2\times 4} \sslash GL_2, \]
    其中不稳定轨迹由秩严格小于 \(2\) 的矩阵组成。
\end{exm}

\begin{rmk}
    当 \(G\) 阿贝尔时，对一般稳定性条件稳定轨迹上的条件总是满足的，但对非阿贝尔的 \(G\) 这并不总是可能。
\end{rmk}

\begin{rmk}
    对一般的光滑射影簇，原点未必是 lci 的，因此不清楚如何以仍能定义虚圈的方式定义拟映射。
\end{rmk}

\begin{defn}
    到 \(V \sslash G\) 的一个\textit{预稳定拟映射}是映射 \(u \colon C \to [V/G]\)，使得基轨迹
    \[ u^{-1}([V^{\mr{un}}(\theta) / G]) \]
    有限并远离特殊点。由定义，\(u\) 由一个主 \(G\)-丛 \(P \to C\) 连同一个等变映射 \(P \to V\) 组成，这等价于相伴丛 \(P \times_G V\) 的一个截面。
\end{defn}

\begin{exm}
    对五次的情形，这恰好恢复了一个线丛 \(\mc{L} \in \Pic C\)（等价于一个 \(\C^{\times}\)-torsor），再配上 \(\on{Tot}(\mc{L}^{\oplus 5})\) 的一个满足定义方程的截面。
\end{exm}

为了定义次数，我们可以考虑 \([V/G]\) 上的线丛 \(\mc{L}_{\theta}\)。然后可以定义 \(C' \subset C\) 在 \(u\) 下的次数为
\[ \deg (u^* \mc{L}_{\theta}|_{C'}). \]

\begin{exm}
    在 \(G(2,4)\) 的例子中，注意主 \(\on{GL}_2\)-丛等价于 \(C\) 上的秩 \(2\) 向量丛 \(\mc{E}\)。因此 \(P \times_G V\) 等价于 \(\mc{E}^{\oplus 4}\) 的数据，从而到格拉斯曼簇的拟映射等价于稳定商。此时我们要找的是在一般点处满射的
    \[ \mc{O}^{\oplus 4} \to \mc{E}, \]
    而由于 \(\theta = \det\)，我们看到 \(u^* \mc{L}_{\theta} = \det \mc{E}\)。
\end{exm}

\begin{exm}
    考虑 \(V = \C^{N+1}\)，注意 \(\P^N = V \sslash \C^{\times}\)。若取 \(\theta_1 \colon z \mapsto z\)，则有
    \[ V^{\on{un}} = V((x_i)). \]
    然而若取 \(\theta \colon z \mapsto z^2\)，则有
    \[ V^{\on{un}} = V((x_i x_j)). \]
\end{exm}

由此可见，概形结构对特征标的选取十分敏感，因此我们不能简单地按朴素方式定义基点的长度。注意
\[ V^{\on{un}} = V\ab(\ab< \bigoplus_{k \geq 1} R_k>), \qquad R_k = H^0([V/G], L_{\theta}^{\otimes k}). \]
特别地，我们必须定义
\[ \on{len}(x) \coloneqq \min_{\substack{f \in R_k \\ k \geq 1}} \frac{\on{ord}_x u^*(f)}{k}. \]

\begin{exm}
    设 \(X = \mr{pt} = \C \sslash_{\theta} \C^{\times}\)。则预稳定拟映射就是
    \[ (C, \mc{L}, \varphi \in H^0(\mc{L})), \]
    其中 \(\varphi^{-1}(0)\) 有限并远离特殊点。此外，若 \(\deg \mc{L} = 1\)，我们恰好恢复了 \(C^{\mr{sm}}\) 中一个新标记点的数据。
\end{exm}

\begin{exm}
    现在考虑 \(X = \mr{pt} = \C^n \sslash (\C^{\times})^n\)。再设次数为 \((1, \ldots, 1)\)。现在取
    \[ \theta = (a_1, \ldots, a_n). \]
    于是注意
    \[ \ab<R_1> = \ab<(x_1^{a_1} \cdots x_n^{a_n})>. \]
    由于这些都位于次数 \(1\)，我们看到对光滑点 \(x \in C^{\mr{sm}}\) 有
    \[ \on{len}_x = \sum \on{ord}_x (\varphi_i^{a_i}) = \sum a_i \on{ord}_x \varphi_i. \]
    若令 \(\ab\{p_i\} = \varphi_i^{-1}(0)\)，则条件 \(\ep \cdot \on{len}(x) \leq 1\) 等价于要求
    \[ \sum_{x = p_i} \ep a_i \leq 1. \]
    次数条件现在变为
    \[ \ep \sum_i a_i + \deg \omega_C > 0. \]
    等价地，这说明
    \[ \omega_C(\ep a_1 p_1 + \cdots + \ep a_n p_n) > 0, \]
    于是我们得到带权标记点曲线，其权重为 \((\ep a_1, \ldots, \ep a_n)\)。
\end{exm}

\subsection{主空间}%
\label{sub:Master space}

证明使用了虚局部化~\cite{GP1999}；Hassett\textquotesingle s 加权模空间~\cite{Hassett2003} 为该构造提供了一个模型。

我们的方法是在插值于两个模空间之间的主空间上作虚局部化。考虑 \(\P^n\) 上由下式给出的 \(\C^{\times}\)-作用
\[ t \cdot [x_0 : \cdots : x_n] = [t x_0, x_1, \ldots, x_n]. \]
不动点轨迹由一个 \(x_0 = 0\) 的 \(\P^{n-1}\) 与一个点 \([1, 0, \ldots, 0]\) 组成。于是我们可以通过把计算转移到另一个不动点轨迹来算出
\[ \int_{\P^{n-1}} H^{n-1} = 1 \]
。特别地，我们提升到 \(H \in H_{\C^{\times}}^* (\P^n) = \Q[H,z] / (H-z)H^n\) 并计算
\begin{align*}
    \int_{\P^n} H^{n-1} &= \int_{\P^{n-1}} \frac{H^{n-1}}{-z+H} + \int_{\mr{pt}} \frac{z^{n-1}}{z^n}.
\end{align*}
取留数，我们看到
\begin{align*}
    0 &= \on{Res}_{z = 0}\int_{\P^{n-1}} \frac{H^{n-1}}{-z+H} + 1 \\
    &= 1 - \int_{\P^{n-1}}H^{n-1},
\end{align*}
这就恢复了所要的积分。我们在此用到的关键事实是射影空间是真（proper）的。

\begin{exm}[第二个玩具例子]
    我们再考虑第二个玩具例子。回忆 Hassett 的加权标记点曲线模空间
    \[ (C, x_1, \ldots, x_n) \]
    其中 \(C\) 是结点曲线，\(x_i \in C^{\mr{sm}}\) 可以碰撞。碰撞受到如下约束：每个标记点被赋予一个权重 \(\delta_i \in \Q \cap (0,1]\)，并要求满足以下条件：
    \begin{enumerate}
        \item 对所有 \(x \in C^{\mr{sm}}\)，有
            \[ \sum_{x = x_i} \delta_i \leq 1; \]
        \item 我们要求稳定性条件 \(\omega_C(\delta_1 x_1 + \cdots + \delta_n x_n) > 0\)。例如，若有一个有理尾，则 \(\sum_{x_i \in C'} \delta_i > 1\)。
    \end{enumerate}
\end{exm}

\begin{exm}[另一个玩具例子]
    我们的目标是再次计算
    \[ \int_{\bar{M}_{0,4}} \psi_1 = 1. \]
    我们的第一步是改变稳定性条件。若取
    \[ \delta_+ = \ab(\frac{2}{3}, \frac{2}{3}, \frac{2}{3}, \frac{2}{3}), \]
    可见稳定性条件没有变化。然而若取
    \[ \delta_- = \ab(\frac{2}{3}, \frac{2}{3}, \frac{2}{3}, \frac{1}{3}), \]
    模空间并未改变，但泛曲线实际上变成 \(\P^1 \times \P^1\)，其中截面 \(x_4\) 是对角线。这告诉我们
    \[ \int_{\bar{M}_{0,4}^{\delta_-}} \psi_1 = 0, \]
    我们将用它在 \(\delta_+\) 的室中计算。我们可以尝试通过既允许带有 \(x_i, x_4\) 的有理尾、又允许 \(x_1 = x_4\) 来构成一个主空间，但若这样做，我们最终会得到一个带有 \(x_1 = x_4\) 的有理尾的半稳定点。我们希望通过从主空间取商来得到它，因此我们给模空间添加一个切向量的数据。于是，定义 \(\delta_{\pm}\) 主空间由如下形式的对象组成
    \[ \xi = (C, x_1, \ldots, x_4, v \in T_{x_4} C \cup \infty) \]
    其中 \(C\) 是结点曲线，\(x_i\) 是光滑点，\(x_1, x_2, x_3\) 互不相同，
    \[ \sum_{x_i \in C'} \frac{2}{3} \delta_i > 1 \]
    对所有有理尾 \(C'\) 成立，若 \(v = 0\)，则 \(x_4\) 与其他 \(x_i\) 不同，而当 \(v = \infty\) 时，不存在 \(x_4 \neq x_i\)（\(i = 1,2,3\)）的有理尾。

    现在我们通过缩放 \(v\) 给这个主空间一个 \(\C^{\times}\)-作用。不动点轨迹有以下几种：
    \begin{enumerate}
        \item 当 \(v = 0\) 时，得到 \(\bar{M}_{0,4}^{\delta_+}\)；
        \item 当 \(v = \infty\) 时，得到 \(\bar{M}_{0,4}^{\delta_-}\)；
        \item 还有一些不动点轨迹，其中 \(v \neq 0,\infty\) 且有一条满足 \(x_i = x_4\) 的有理尾。
    \end{enumerate}

    若考虑一个一般点（四个标记点互不相同），它沿 \(\C^{\times}\)-作用直接走到自身。然而若我们从一条包含 \(x_1\) 与 \(x_4\) 的有理尾出发，则先走到一条 \(v\) 一般但 \(x_1 = x_4\) 的有理尾，若把结点光滑化，就得到一条 \(x_1 = x_4\) 的单条曲线，而我们允许 \(v \to \infty\)，这是另一个不动点。现在我们可以积分
    \begin{align*}
        \int_{\ms{Master}} \psi_1 = \int_{\bar{M}_{0,4}^{\delta_+}} \frac{\psi_1}{-z-?} + \int_{\bar{M}_{0,4}^{\delta_-}} \frac{\psi_1}{z+?} + \on{Cont}(12|34) + \on{Cont}(13|24) + \on{Cont}(14|23).
    \end{align*}
    在 \(z = 0\) 处取留数，我们得到
    \begin{align*}
        0 &= - \int_{\bar{M}_{0,4}^{\delta_+}} \psi_1 + 0 + 0 + 0 + \int_{\mr{pt}} \frac{(-z)}{(\pm z)(\pm z)},
    \end{align*}
    其中最后一项分母中的 \(\pm z\) 是结点光滑化参数，于是我们恢复了所要的积分。
\end{exm}

有了所有这些玩具例子，我们终于可以讨论拟映射主空间了。我们在 \(\frac{1}{d_0}\) 处的墙上的第一步是同时允许次数 \(d_0\) 的有理尾与长度 \(d_0\) 的基点。然而，我们可以考虑带有如下形式拟映射之有理尾的曲线
\[ [a_0 u^{d_0}, \ldots, a_n u^{d_0}], \]
它们有无穷多个自同构。然而，我们不能在基点上附加切向量（因为形变会把它们消去），而且也存在多于一个次数 \(d_0\) 并带有长度 \(d_0\) 基点的情形（此时自同构群包含 \(\C^{\times}\)）。

\begin{rmk}
    有一个关键观察。假设我们有两条次数 \(d_0\) 的有理尾，并称两个结点为 \(p_1\) 与 \(p_2\)。现在定义光滑化参数的空间
    \[ \C \cong \Theta_i = T_{p_i} E_i \otimes T_{p_i} C_g, \]
    其中 \(E_i\) 是有理尾，\(C_g\) 是曲线的其余部分。我们考虑由次数标号的加权半稳定曲线构成的模空间
    \[ \mf{M}_{g,n,d}^{\mr{wt}} \xrightarrow{\text{\'et}} \mf{M}_{g,n} \]
    注意该态射高度非分离。我们考虑开子叠 \(\mf{M}_{g,n,d}^{\mr{wt}, \ep_0}\)，其中的曲线没有次数严格小于 \(d_0\) 的尾。然后我们可以考虑 Cartier 除子 \(\mc{Z}\)，即存在至少一条次数 \(d_0\) 的尾的地方。特别地，对所关注点处的两个分支 \(\mc{Z}_1\) 与 \(\mc{Z}_2\)，有
    \[ \Theta_i = N_{\mc{Z}_i} |_{\xi} \]
    。我们还看到 \(\Theta_1 \otimes \Theta_2\ = \mc{O}(\mc{Z})|_{\xi}\)，因此可以添加数据
    \[ v \in \ab(\bigotimes_i \Theta_i)^{\vee} \cup \ab\{\infty\} \]
    以及数据
    \[ w \in \P(\Theta_1 \oplus\Theta_2). \]
\end{rmk}

实际构造如下进行。考虑 \(\mc{Z} \subset \mf{M} = \mf{M}_{g,d}^{\mr{wt}, \ep_0}\)。定义 \(\tilde{\mf{M}} \to \cdots \to \mf{M}\) 为沿 \(\mc{Z}\) 的各层（strata）的一系列爆破。于是主空间在概形 \(S\) 上的一个对象（带校准尾的拟映射）由如下数据给出：一个拟映射
\begin{equation*}
\begin{tikzcd}
    C \ar{d} \ar[dashrightarrow]{r}{f} & X \\
    S
\end{tikzcd}
\end{equation*}
它诱导态射 \(p \colon S \to \mf{M}_{g,d}^{\mr{wt}, \ep_0}\)，一个提升的选取
\begin{equation*}
\begin{tikzcd}
    & \tilde{\mf{M}} \ar{d}{\tau} \\
    S \ar{r}{p} \ar[dashrightarrow]{ur}{w} & \mf{M}_{g,d}^{\mr{wt}, \ep_0},
\end{tikzcd}
\end{equation*}
一个线丛 \(\mc{N} \in \Pic S\)，以及 \(v = (v_1, v_2)\)，其中 \(v_1 \in H^0(w^* \mc{O}(-\mc{Z}))\)，\(v_2 \in H^0(\mc{N})\)，且 \((v_1, v_2)\) 不全为零。我们还禁止长度严格大于 \(d_0\) 的基点，并要求如下稳定性
\begin{enumerate}
    \item 当 \(v = 0\) 时，得到 \(\ep_+\)-稳定性；
    \item 当 \(v = \infty\) 时，得到 \(\ep_-\)-稳定性；
    \item 任何常尾（次数 \(d_0\) 且带有长度 \(d_0\) 的基点）都是纠缠尾。这里我们考虑最后一次爆破使得 \(w\) 映到例外轨迹 \(\bigcap_{i \in I} \tilde{\mc{Z}}_i\)，此时对 \(i \in I\)，\(E_i\) 都是纠缠尾；
    \item 若存在次数 \(d_0\) 的尾，则所有长度 \(d_0\) 的基点都在次数 \(d_0\) 的尾上。
\end{enumerate}







